\documentclass[../../../include/open-logic-section]{subfiles}

\begin{document}

\olfileid{sth}{choice}{justifications}
\olsection{اعتبارات داخلية بشأن الاختيار}

السؤال الأوسع إذن هو هل يمكن تسويغ الترتيب الجيد، أو الاختيار، أو
حتى قابلية جميع المجموعات للمقارنة من حيث حجمها---ولا يهم أيها.

وفيما يأتي محاولة لتسويغ \emph{داخلي}. قدمنا في
\olref[z][story]{sec} عدة مبادئ عن الهرم. ويجدر بنا إعادة ذكر أحدها:
\begin{enumerate}
	\item[] \stagesacc. لكل مرحلة~$S$، ولأي مجموعات تكونت
	\emph{قبل} المرحلة~$S$: تتكون في المرحلة~$S$ مجموعة عناصرها هي تلك
	المجموعات بالضبط. ولا يتكون شيء آخر في المرحلة~$S$.
\end{enumerate}
والواقع أن كثيرًا من المؤلفين اقترحوا إمكان تسويغ مسلّمة الاختيار
بواسطة (شيء يشبه) هذا المبدأ. وسنقدم بإيجاز شرحًا لذلك النهج.

سنبدأ بنتيجة صغيرة بسيطة، تقدم مكافئًا \emph{آخر أيضًا} للاختيار:

\begin{thm}[في $\ZF$]\ollabel{choiceset}
الاختيار مكافئ للمبدأ الآتي. إذا كانت !!{element}s~$A$ منفصلة وغير
خالية، فيوجد~$C$ بحيث تكون~$C \cap x$ مجموعة أحادية لكل~$x \in A$.
(ونسمي~$C$ كهذه {مجموعة اختيار} لـ$A$.)
\end{thm}

برهان هذه النتيجة مباشر، ونتركه تمرينًا للقارئ.

\begin{prob}
أثبت~\olref[sth][choice][justifications]{choiceset}. وإذا تعثرت، تجد
برهانًا في~\cite[pp.~242--3]{Potter2004}.
\end{prob}

النقطة الجوهرية هي أن مجموعة اختيار لـ$A$ ليست إلا مدى دالة اختيار
لـ$A$. ولذلك، كي نسوّغ الاختيار، يمكننا ببساطة أن نحاول تسويغ صيغته
المكافئة، من حيث وجود مجموعات الاختيار. وسنحاول الآن فعل ذلك بالضبط.

لتكن !!{element}s~$A$ منفصلة وغير خالية. بموجب~\stageshier{}
(انظر~\olref[z][story]{sec})، تتكون~$A$ في مرحلة ما~$S$. لاحظ أن جميع
!!{element}s~$\bigcup A$ متاحة قبل المرحلة~$S$. والآن، بموجب
\stagesacc{}، تتكون، \emph{لأي} مجموعات تكونت قبل~$S$، مجموعة عناصرها
هي تلك المجموعات بالضبط. وبعبارة أخرى: توجد في~$S$ كل \emph{مجموعة
ممكنة} من المجموعات المتاحة سابقًا. لكن من \emph{الممكن} قطعًا
اختيار أشياء يمكن تكوين مجموعة اختيار لـ$A$ منها؛ فما ذلك إلا مجموعة
جزئية محددة جدًا من~$\bigcup A$. ومن ثم توجد مجموعة اختيار كهذه، كما
هو مطلوب.

حسنًا، تلك محاولة \emph{بالغة} السرعة لتقديم تسويغ للاختيار على أسس
داخلية. لكن لمواصلة بحث هذه الفكرة، ينبغي لك قراءة معالجة بوتر
الأنيقة لها (\citeyear[\S14.8]{Potter2004}).

\end{document}
